\chapter{SCINE Integration}
\label{ch:scine-integration}

SCINE Chemoton explores a chemical reaction network in cycles of
elementary-step proposals and refinements, and the daemon SCINE Puffin
executes the individual jobs \cite{unsleber2022chemoton,weymuth2024scine,bensberg2025puffin}.
Neither program accepts a calculator object: every job resolves its calculator
from the \code{method_family} string of the database model, at a single point
in \file{scine_puffin/utilities/scine_helper.py}.  That single point is what
\jns{} intercepts.  This chapter describes the calculator bridge that lets an
ASE calculator answer SCINE's \code{Calculator} interface
(Section~\ref{sec:bridge}), the process-level patch that routes the declared
method families to it (Section~\ref{sec:injection}), the readuct-level route
kept for comparison (Section~\ref{sec:readuct-patch}), and the deployment and
verification procedure (Section~\ref{sec:scine-verify}).

\section{Why a bridge is needed}
\label{sec:why-bridge}

Puffin's \code{SettingsManager.prepare_readuct_task} is the only place where
a readuct-type job turns a method name into a calculator:

\begin{excerpt}
scine_puffin/utilities/scine_helper.py
    SettingsManager.prepare_readuct_task(...)
        system = utilities.core.load_system_into_calculator(
            "system.xyz", model.method_family, **calculator_settings)
\end{excerpt}

The engine never sees a calculator; it sees a string.  The same constructor
path also calls \code{get_available_settings(method_family, program)} (the
C++ side answers this call with \code{getCalculator(...)->settings()}), so
leaving that function unpatched makes \code{SettingsManager} fail at
construction, before any calculation can run.

\jns{} supplies two building blocks for this point: the bridge
(\code{jnuscout.scine.bridge.ASECalculatorBridge}), a Python subclass of the
SCINE \code{Calculator} base class that wraps any ASE calculator, and the
injection patch (\code{jnuscout.scine.puffin_patch}), which replaces the
resolution functions in-process.

\section{The calculator bridge}
\label{sec:bridge}

\subsection{The interface contract}

A Python calculator in SCINE Core is driven through double dispatch: public
methods of the C++ base class call back into \code{_xxx_impl} methods on the
Python subclass, and a missing override aborts with

\begin{transcript}
Missing overload of '_xxx_impl' in Python Calculator derivative.
\end{transcript}

\code{ASECalculatorBridge} implements all sixteen of them
(Table~\ref{tab:bridge-impl}).

\begin{table}[htbp]
\centering
\caption{The sixteen \code{_*_impl} methods of the SCINE Calculator
contract, in the order of the dispatch list in
\file{src/jnuscout/scine/bridge.py}.}
\label{tab:bridge-impl}
\begin{tabularx}{\textwidth}{@{}lX@{}}
\toprule
Method & Role \\
\midrule
\code{_calculate_impl} & runs the ASE backend, writes the results object and returns it \\
\code{_clone_impl} & builds a fully isolated copy (settings and structure) \\
\code{_results_impl} & returns the results object; the C++ side keeps the reference \\
\code{_get_structure_impl} & returns the \code{AtomCollection} \\
\code{_set_structure_impl} & stores an \code{AtomCollection} \\
\code{_get_positions_impl} & returns the structure positions \\
\code{_modify_positions_impl} & writes positions into the structure \\
\code{_settings_impl} & returns the descriptor-carrying \code{Settings} object \\
\code{_set_required_properties_impl} & stores the required-property list \\
\code{_get_required_properties_impl} & returns the required-property list \\
\code{_possible_properties_impl} & declares the property superset the bridge provides \\
\code{_supports_method_family_impl} & accepts every method family \\
\code{_get_state_impl} & returns an empty SCINE state \\
\code{_load_state_impl} & accepts a state without storing it \\
\code{_name_impl} & returns the calculator name (diagnostic) \\
\code{_allows_python_gil_release_impl} & returns \code{False}: ASE and torch callbacks need the GIL \\
\bottomrule
\end{tabularx}
\end{table}

All quantities cross the boundary in SCINE's atomic units (Hartree, Bohr),
while the ASE side uses eV and \AA, so the bridge converts on both axes:

\begin{equation}
  E[\mathrm{E_h}] = \frac{E[\mathrm{eV}]}{27.211386245}, \qquad
  g[\mathrm{E_h}/\mathrm{bohr}] = -\,\frac{F[\mathrm{eV}/\text{\AA}]}
      {1.8897259886 \times 27.211386245}.
  \label{eq:bridge-units}
\end{equation}

Note the sign convention in the second relation: a force is the negative
gradient.

\subsection{Return the results object}
\label{sec:contract-return}

The C++ trampoline for \code{calculate} is declared as

\begin{excerpt}
PYBIND11_OVERRIDE_PURE_NAME(const Scine::Utils::Results&, Scine::Core::Calculator,
                            "_calculate_impl", calculate, description)
\end{excerpt}

so the value returned by \code{_calculate_impl} must convert to a
\code{const Results\&}.  Returning \code{None} makes the pybind11 conversion
raise a \code{TypeError}; \code{CalculationRoutines::calculateWithCatch}
swallows any exception with \code{catch (...)}, prints only

\begin{transcript}
Aborting optimization due to failed calculation
\end{transcript}

and rethrows.  The observable symptom is a geometry optimisation that stops
after a single evaluation, with a message that points nowhere near the cause.
The bridge therefore ends \code{_calculate_impl} with \code{return r}, and
sets \code{successful_calculation} as the last write to the object, because
the C++ side uses that flag to decide whether the step succeeded.

\subsection{Update the results object in place}
\label{sec:contract-inplace}

readuct reads the results object through \code{_results_impl} before the
calculation and keeps the reference (three reads ahead of a single-point
task).  Re-assigning \code{self._results = su.Results()} inside
\code{_calculate_impl} leaves the held reference pointing at an object whose
\code{successful_calculation} is never set, with the same external symptom as
above.  The bridge writes the fields of the existing object instead of
replacing it.

\subsection{Clones must be fully isolated}
\label{sec:contract-clone}

readuct clones the calculator during argument handling and runs the
calculation on the clone.  Puffin additionally derives spin-shifted systems by
cloning a calculator and writing a new multiplicity into the clone:

\begin{excerpt}
shifted = systems[name].clone()
shifted.settings["spin_multiplicity"] = <shifted multiplicity>
\end{excerpt}

The propensity iterator is lazy and re-reads the multiplicity of the original
object on every round, so a clone that shares the settings object pollutes the
original with the last written shift value, and the job fails with

\begin{transcript}
Could not find system nt_multiplicity_shift_-2
\end{transcript}

The bridge therefore builds a fresh descriptor-carrying \code{Settings} object
on clone and copies keys one by one, and it rebuilds the structure as a new
\code{AtomCollection} from elements and positions.  Sharing the structure
would make \code{_modify_positions_impl} on the clone move the original
geometry as well, because the optimiser operates on clones throughout.

\subsection{A declared superset of properties and settings}
\label{sec:superset}

The engine fails the whole job when it requires a property the calculator did
not declare, so the bridge over-declares.  \code{_possible_properties_impl}
declares energy, gradients, atomic charges, the bond-order matrix and the
Hessian.  The error messages that motivated the additions are

\begin{transcript}
Charges required, but chosen calculator does not provide them
Bond orders required, but chosen calculator does not provide them
Calculator is missing a Hessian property
\end{transcript}

Each declared property is produced when it is requested:

\begin{itemize}
  \item energy and gradients come from the ASE backend, converted as in
        equation~\eqref{eq:bridge-units};
  \item atomic charges are a placeholder of zeros: the MACE and AIMNet2
        backends do not expose charges, and reactive-site detection in the
        AFIR search relies mainly on bonding topology;
  \item bond orders are detected on the SCINE-unit structure with SCINE's own
        distance criterion (\code{su.BondDetector.detect_bonds}); the
        behaviour is controlled by \code{NNP_AFIR_BOND_ORDERS} (on by
        default);
  \item the Hessian is computed by central differences of the analytic
        gradients with a step of $1\times10^{-2}$~Bohr, which makes the unit
        naturally $\mathrm{E_h}/\mathrm{bohr}^2$; the cost is two gradient
        evaluations per Cartesian coordinate ($6N$ for $N$ atoms), and it is
        only evaluated when the required-property list contains it.
\end{itemize}

The same reasoning applies to settings.  An empty descriptor collection makes
readuct abort while writing \code{self_consistence_criterion} inside
optimisation or scan tasks,

\begin{transcript}
Setting 'self_consistence_criterion' not implemented in the applied calculator
\end{transcript}

and Puffin writes program-layout keys (\code{base_working_directory},
\code{output_path}, \dots) that the official C++ calculators declare
themselves; an undeclared key surfaces as

\begin{transcript}
There is no matching key 'X' in these Settings!
\end{transcript}

\code{bridge_descriptors()} therefore declares the full standard key set for
the bridge (27 keys).  Unknown keys that the engine still sends are skipped
and counted by the patch rather than rejected; with
\code{NNP_AFIR_PUFFIN_VERBOSE=1} the skipped keys are printed.  The
asymmetry is deliberate: over-declaring costs nothing, under-declaring fails
the job before it starts.

\begin{refusalbox}[What the bridge does not provide]
The bridge is an adapter, not a model.  It does not compute atomic charges
(the charge field is a placeholder), it does not persist calculator state
(\code{_get_state_impl} returns an empty state and \code{_load_state_impl}
stores nothing), and it does not release the GIL during a calculation.  It
also does not judge: whether a prediction is trustworthy is decided by the
uncertainty layer, not here.
\end{refusalbox}

\section{Injection into Puffin and readuct}
\label{sec:injection}

\subsection{The patch}
\label{sec:patch}

\code{puffin_patch.install_from_env()} replaces the calculator-resolution
functions on \code{scine_utilities.core} for the declared method families:

\begin{itemize}
  \item \code{load_system_into_calculator} returns a bridge whose structure
        is the system the engine passed in.  An \code{AtomCollection} is used
        as-is with no unit conversion (SCINE works in Bohr throughout); a file
        path is read from disk and converted to Bohr;
  \item \code{get_available_settings} returns a descriptor-carrying
        \code{Settings} object equivalent to \code{getCalculator(...)->settings()};
  \item \code{has_calculator} and \code{get_calculator} answer the companion
        queries for the same families.
\end{itemize}

Every other method family falls back to the original function, so
Sparrow/DFTB3, PM6 and DFT jobs keep working unchanged.  One calling
convention needs care: \code{SettingsManager.calculator_settings} itself
contains a \code{method_family} key, so the engine calls

\begin{excerpt}
load_system_into_calculator(xyz, method_family, **settings)
\end{excerpt}

with the family both positionally and as a keyword.  The wrapper accepts this
leniently and de-duplicates; a strict signature would fail on the first job.

\subsection{Environment variables}
\label{sec:scine-env}

\begin{table}[htbp]
\centering
\caption{Environment variables of the injection patch.  The names are kept
from earlier deployments and are part of the interface.}
\label{tab:scine-env}
\begin{tabularx}{\textwidth}{@{}lXl@{}}
\toprule
Variable & Meaning & Default \\
\midrule
\code{NNP_AFIR_PUFFIN} & enable the patch (called from \code{sitecustomize.py}) & off \\
\code{NNP_AFIR_METHOD_FAMILY} & family, or comma-separated families, to intercept & \code{nnp} \\
\code{NNP_AFIR_BACKEND} & \code{emt}, \code{mace} or \code{aimnet2} & \code{mace} \\
\code{NNP_AFIR_MODEL} & path to the MACE model file & unset \\
\code{NNP_AFIR_DEVICE} & torch device for MACE & \code{cuda} \\
\code{NNP_AFIR_BOND_ORDERS} & attach the bond-order matrix to the results & on \\
\code{NNP_AFIR_PUFFIN_VERBOSE} & print hit messages and skipped settings & off \\
\code{NNP_AFIR_BRIDGE_VERBOSE} & print bridge-level diagnostics & off \\
\bottomrule
\end{tabularx}
\end{table}

The \code{emt} backend builds an ASE EMT calculator: it needs no model weights
and no torch, which makes it the lightest end-to-end smoke test of the
plumbing.  The \code{aimnet2} backend is not wired up yet; selecting it raises
\code{NotImplementedError}.  For production use, point
\code{NNP_AFIR_MODEL} at a fine-tuned MACE model.

\subsection{Deployment: activate before Puffin starts}
\label{sec:scine-deploy}

The patch must be installed in every process that may execute a job, including
forked workers, and before Puffin's main routine starts.  The supported
mechanism is a \code{sitecustomize.py} on \code{PYTHONPATH}: Python imports it
automatically at interpreter start-up, so the daemon and everything it forks
inherit the patch.

\begin{codeblock}
# sitecustomize.py -- lives in a directory placed on PYTHONPATH
import jnuscout.scine as scine

scine.install_from_env()
\end{codeblock}

\begin{transcript}
$ export PYTHONPATH=/opt/jnu-scout/src:/opt/jnu-scout-deploy
$ export NNP_AFIR_PUFFIN=1
$ export NNP_AFIR_METHOD_FAMILY=nnp
$ export NNP_AFIR_BACKEND=mace
$ export NNP_AFIR_MODEL=/opt/models/mace_al_v3f.model
$ puffin start -c puffin.yaml
\end{transcript}

The installation is idempotent, and \code{scine.uninstall()} restores the
original functions.  \code{install_from_env()} is a no-op unless
\code{NNP_AFIR_PUFFIN} is set to a non-empty, non-zero value.

\section{The readuct-level patch (kept for comparison)}
\label{sec:readuct-patch}

\code{jnuscout.scine.readuct_patch} predates the full bridge and takes a
different route.  It wraps every \code{readuct.run_*_task} function, captures
the calculator clones that readuct makes during argument handling, and, when a
task raises the specific return-value conversion \code{TypeError}
(\code{incompatible function arguments}) after a successful calculation,
rebuilds the returned dictionary from the captured clone.
\code{install_puffin_patch(bridge_cls, factory, sentinel="nnp")} additionally
replaces \code{ScineHelper.prepare_readuct_task} for the sentinel method
family.

This module is a workaround for a binding defect, not a fix, and it is kept in
the package for comparison only:

\begin{itemize}
  \item it matches one specific \code{TypeError} string and applies only to
        task functions that return a calculator dictionary;
  \item if a task fails for a different reason, no clone is captured and the
        wrapper re-raises; set \code{NNP_PATCH_VERBOSE=1} to print a notice
        whenever it does bypass a conversion.
\end{itemize}

\begin{notebox}[Recommended route]
Use \code{puffin_patch.install_from_env()}.  With the interface contract of
Section~\ref{sec:bridge} satisfied, the calculation path needs no return-value
workaround; \code{readuct_patch} exists to document and reproduce the
earlier failure mode.
\end{notebox}

A third, factory-level variant lives in \code{jnuscout.scine.inject}
(\code{apply_patch()}).  It routes the families \code{("DFT", "GFN2-xTB")} to
\code{ScineORCACalculator}, which implements only five of the sixteen
\code{_*_impl} methods and has no \code{_clone_impl}; that chain is
incomplete by itself and is kept for comparison as well.

\section{Verifying the integration}
\label{sec:scine-verify}

A minimal check exercises the whole path -- database model, Puffin job
resolution, bridge, backend -- with the lightest possible backend:

\begin{enumerate}
  \item start MongoDB and Puffin with the patch active and with
        \code{NNP_AFIR_BACKEND=emt};
  \item build a small database whose model carries
        \code{method_family='nnp'} and submit one exploration cycle;
        \script{setup_h2co_nnp_db} does exactly this for the H$_2$CO test
        system;
  \item inspect the job status in the database.
\end{enumerate}

\begin{readbox}[How to read the result]
A job that reaches \code{Status.COMPLETE} has executed the calculation and the
job's post-processing (structure and energy bookkeeping, bond orders, graph
updates) without the bridge raising.  If a job fails instead, re-run with
\code{NNP_AFIR_PUFFIN_VERBOSE=1} and \code{NNP_AFIR_BRIDGE_VERBOSE=1}:
the first prints which resolution calls hit the patch, the second prints
bridge-level diagnostics.  The message ``Aborting optimization due to failed
calculation'' is generic by construction -- it means only that
\code{calculate()} raised.  Check the return-value contract of
Section~\ref{sec:contract-return} first.
\end{readbox}

Two silent failure modes are worth remembering on a new deployment: a job that
never starts is often a missing settings descriptor, because the engine
rejects an unknown key before the first calculation; a job that stops after a
single evaluation is usually a violated return-value or in-place-results
contract.
